\section{Analisi del Problema}
L'implementazione di \emph{Poool} nasce dall'unione dei due sketch: \texttt{sketch01}, orientato alla logica sequenziale, mentre \texttt{sketch02}, orientato all'interazione asincrona in architettura MVC. 
Durante questa integrazione sono emerse alcune criticità:

\subsection{Input asincroni e reattività dell'EDT}
\label{sec:analisi_edt}
\textbf{L'Event Dispatch Thread (EDT)} di Swing deve restare reattivo: ogni operazione bloccante sull'EDT produce freeze percepibili della GUI. È necessario l'utilizzo di un thread separato, per garantire sia l'aggiornamento della View, sia la gestione degli input asincroni.

Il problema è risolto tramite l'architettura MVC con observer: implementando un \textbf{Controller} in cui i comandi utente devono essere trasferiti dalla View ad esso senza perdere eventi e senza introdurre attese indefinite. 
Nello \texttt{sketch02} viene utilizzato un \textbf{AutonomousUpdater}, un thread dedicato ad aggiornare periodicamente lo stato di un Counter. Ho trovato similarità tra questo pattern e la necessità di un loop di gioco che avanzi a cadenza regolare, indipendentemente dagli input.

\subsection{Parallelizzazione delle collisioni}
La fase più costosa del motore è la risoluzione delle collisioni tra palline, basata su confronti a coppie. Con $n$ palline, il costo cresce quadraticamente come $ \mathcal{O}(n^2) $.
Negli scenari ad alta densità (la configurazione \textbf{massive}: 4500 palline), un calcolo sequenziale non riesce a mantenere la stessa fluidità dei casi con meno palline, rendendo necessaria una parallelizzazione mirata del solo sottoproblema delle collisioni.
Un primo approccio efficace alla parallelizzazione consiste nella partizione della lista di palline, garantendo uniformità al calcolo delle collisioni tra i thread worker simultanei.

\subsubsection*{Coordinamento e Coerenza}
Parallelizzare le collisioni assegnando blocchi di palline a thread worker è solo una parte della soluzione; è necessario anche coordinare i thread per garantire che il game loop resti coerente.
Affinché il rendering mostri uno stato complessivo corretto e veritiero risulta fondamentale imporre un punto di ritrovo o attesa condiviso (sfruttando il concetto di \textbf{Barrier}). Questo vincolo assicura l'integrità del game loop.

\subsubsection*{Rischi di Deadlock nelle Collisioni}
La risoluzione concorrente su coppie condivise di palline richiede lock multipli. Senza una politica deterministica di acquisizione, due thread potrebbero ottenere i lock in ordine inverso, causando attesa circolare e quindi deadlock.

Per prevenire questa condizione viene applicato un lock ordering globale sulle palline: i monitor sono acquisiti sempre nello stesso ordine logico. La validità della scelta è stata poi verificata formalmente con JPF \ref{sec:verifica_formale}.